\documentclass[12pt]{article}

\usepackage{amsmath, amsfonts, amssymb, enumerate}
\usepackage[margin=1cm]{geometry}

\begin{document}

\noindent
{\large\bfseries Number Theory}
\hfill
znz
\quad
\today

\hrule
\medskip

\begin{enumerate}
  \item
        Extend $\deg(0) = -\infty$.

        \[
        \begin{gathered}
          f = qg + r \\
          f(t) = f_{d}t^{d} + \cdots, f_d \neq 0 \\
          q(t) = a_{n}t^{n} + \cdots, a_n \neq 0 \\
          g(t) = b_{m}t^{m} + \cdots, b_m \neq 0 \\
        \end{gathered}
        \]

        For existence:
        \[
        \begin{gathered}
        P(n)\colon \text{for every }f\in F[t]\text{ with }\deg(f)=n, n \in \mathbb{N}, g \neq 0, \exists q, r \colon f=qg + r, \deg(r) < \deg(g) \\
        f = 0 \colon 0 = 0 \cdot g + 0 \\
        \deg(0) = -\infty < \deg(g)
        \end{gathered}
        \]
        \begin{itemize}
          \item base
                \[
                \begin{gathered}
                  P(n) \text{ is true for } n < \deg(g) \\
                  q=0, r=f, f=qg+r, \deg(r) = \deg(f) = n < \deg(g)
                \end{gathered}
                \]
          \item induction step, assume for every $n \leq N \colon P(n)$ is true, prove it for $N + 1$.

                \[
                \begin{gathered}
                  g_{m} \in F, F \text{ is a field }, g_{m}\neq 0 \implies \exists g_{m}^{-1} \\
                  \text{let } a=g_{m}^{-1}f_{N+1} \\
                  \text{let } f'=f - agt^{N+1-m} \implies \deg(f') < \deg(f) \implies
                  f'=q'g + r \implies f=(q' + at^{N+1-m})g + r
                \end{gathered}
                \]
                which concludes the induction.
        \end{itemize}

        For uniqueness assume that $\exists q', r' \colon q \neq q' \text{ or } r \neq r', f=q'g +r'$.
        \[
        \begin{gathered}
          f = qg + r, f = q'g + r' \\
          qg + r - (q'g + r') = 0 \\
          (q - q')g + (r - r') = 0 \\
          \deg(r) < \deg(g), \deg(r') < \deg(g) \implies \deg(r - r') < \deg(g) \implies \deg(r'-r) < \deg(g)
        \end{gathered}
        \]

        Consider two cases:
        \[
        q=q' \implies r=r'
        \]
        contradiction with initial assumption

        \[
        \begin{gathered}
          q \neq q' \implies \deg(q - q') \geq 0 \implies \deg((q - q')g) \geq \deg(g) \\
          (q - q')g = r'-r \implies \deg((q-q')g) = \deg(r'-r)
        \end{gathered}
        \]
        which is a contradiction since $LHS \geq \deg(g), RHS < \deg(g)$, therefore $q'=q, r'=r$.
  \item
        Let's use Euclidean algorithm for finding $\gcd(g, f)$:
        \begin{enumerate}
          \item if $f=0$ result is g
          \item otherwise result is $\gcd(r, g)$ where $f=qg + r$.
        \end{enumerate}

        Such $q, r$ exist (see problem 1), $\deg(r) < \deg(g)$. Therefore steps of Euclidean algorithm will look like:
        \[
        \begin{gathered}
          f=q_0g + r_0 \\
          g=q_1r_0 + r_1 \\
          r_0=q_2r_1 + r_2 \\
          \cdots \\
          r_{k-1}=q_{k+1}r_{k} + r_{k+1}, r_{k+1}=0
        \end{gathered}
        \]

        Since $\deg(r_k)$ decreases with each step this algorithm is going to terminate in a result $d$.

        Let's prove that $d \mid f, d \mid g$. To do that we will do reverse Euclidian algo steps to reconstruct $f, g$ from $d$ using existing coefficients $q_i, r_i$:
        \[
        \begin{gathered}
          r_k=d, r_{k+1}=0 \\
          d \mid r_{i}, d \mid r_{i+1}, r_{i-1}=q_{i+1}r_{i} + r_{i+1} \implies d \mid r_{i-1} \\
          \text{therefore } d \mid r_1, d \mid r_0 \implies d \mid g, d \mid f
        \end{gathered}
        \]

        Now let's consider an arbitrary $c$, such that $c | f, c | g$.
        \[
        \begin{gathered}
          c \mid f, c \mid g \implies c \mid r_0 \\
          c \mid g, c \mid r_0 \implies c \mid r_1 \\
          \cdots \\
          c \mid r_{k}, c \mid r_{k+1} \implies c \mid r_{k+2}
        \end{gathered}
        \]

        Since $\exists k \colon d = r_k$ then $c \mid d$ for arbitrary $c$. Therefore since $d \mid f, d \mid g, d = \gcd(f, g)$.

        Now since $d$ must be monic let's consider $d(t)=d_0 + d_1t + \cdots + d_{n}t^n$. Since $d_{n}\in F, F \text{ --- field} \implies \exists d_{n}^{-1} \implies d'=d_{n}^{-1}d \text{ is monic}$.
  \item
        \begin{enumerate}[a)]
          \item Assume $\gcd(\pi, f)=c$, then $c \mid \pi, c \mid f$. Since $\pi$ is irreducible then $\pi=cd$, and either $c$ or $d$ is a unit. There are two possibilities:
                \begin{itemize}
                  \item $c$ is a unit. But then only monic unit is $1$ and $\gcd(\pi, f)=1$.
                  \item $d$ is a unit, $c$ is an irreducible polynomial. But since $\pi$ is irreducible then $c$ and $\pi$ are from the same equivalence class. But then $\gcd(\pi, f) = \pi$ and $\pi \mid f$ which is a contradiction.
                \end{itemize}
          \item $\pi \mid fg$. Assume $\pi \nmid f \text{ and } \pi \nmid g \implies \gcd(\pi, f) = 1$.
                \[
                \begin{gathered}
                  1 = {\pi}x + fy \\
                  g = {\pi}gx + fgy \\
                  \pi \mid {\pi}gx, \pi \mid fgy \implies \pi \mid g \text{ --- contradiction}
                \end{gathered}
                \]
                therefore $\pi \mid f \text{ or } \pi \mid g$.
          \item define the following algorithm. Let's say $fact(n)$ is a finite product of irreducible polynomials:
                \begin{itemize}
                  \item if $f$ is irreducible $fact(f) = f$
                  \item if $f$ is reducible $\exists h,g \colon f=hg \implies fact(f)=fact(h)fact(g)$.
                  \item $\deg(h) < \deg(f), \deg(g) < \deg(f), \deg \in \mathbb{N}$ and $<$ on $\mathbb{N}$ is well-founded. $NSTEPS(call)$ --- number of steps needed for temination.


                        $P(n)=\forall f \colon \deg(f) = n, NSTEPS(fact(f)) < \infty$

                        $P(1)$ is true since such $f$ are irreducible.

                        Assume $P(n)$ is true for all $n \leq N$. Prove for $N + 1$.

                        For arbitrary $f \colon \deg(f) = N+1, NSTEPS(fact(f))=1 \text{ if it's irreducible or } NSTEPS(fact(f))=NSTEPS(fact(g)) + NSTEPS(fact(h)) + 1 \text{ otheriwse } \implies P(N+1) true$.
                \end{itemize}
          \item consider arbitrary $\pi_i$. If $\pi_i \mid f$ then from 3.b there are two possibilities:
                \begin{itemize}
                  \item $\exists \rho_j \colon \pi_i \mid \rho_j$. This means that $\rho_j$ is either reducible which is a contradiction or $\rho_j=h\pi_i$ where $h \in {F[t]}^{\times}$, which means $\pi_i$ and $\rho_j$ are associates.
                  \item $\pi_i \mid v$, but this is impossible because $\deg(v) < \deg(\pi_i)$.
                \end{itemize}
                To prove 1-to-1 matching of $\pi_i$ and $\rho_j$ and $r=s$ do the following:
                \begin{itemize}
                  \item Remove $u$ from the left hand side by multiplying both sides by $u^{-1}$ and reassigning $v=vu^{-1}$. Since all $\pi_i$ and $\rho_j$ are monic then their products leading coefficients are equal to $1$ which means $v=1$ as well and can be removed.
                  \item Suppose $\pi_1 \cdots \pi_r = \rho_1 \cdots \rho_s$. Then $\exists j \colon \pi_1=\rho_j$.

                        $F[t]$ is an integral domain $\implies$ multiplication is commutative $\implies$ we can reorder and get
                        \[
                        \pi_1 \cdots \pi_r = \rho_j \rho_1 \cdots \rho_{j-1} \rho_{j+1} \cdots \rho_s
                        \]

                        $F[t]$ is an integral domain $\implies 0$ has no non-zero divisors which allows cancellation and reindexing:
                        \[
                        \pi_2 \cdots \pi_r = \rho_2 \cdots \rho_s
                        \]

                        Since number of factors is finite this process repeats until no more factors left on left or right side of the equality.

                  \item Assume $r \neq s$, for example $r < s$ then after this process ends we get
                        \[
                        1=\rho_k \cdots \rho_s
                        \]

                        Whicl is impossible because $\deg(1) = 0, \deg(\rho_k \cdots \rho_s) \geq 1$.

                        Similarly for $r > s$. Contradiction.
                \end{itemize}
        \end{enumerate}

\end{enumerate}

\end{document}
